\chapter[Spike-train renewal and spectra]{Renewal processes, spike-train spectra, and bin limits}
\label{ch:renewal-spectral}
A spike train may be treated as a random counting measure
$n(t)=\sum_j\delta(t-t_j)$, where the Dirac deltas are
\emph{distributions}, not finite-amplitude voltage spikes.
This representation connects point processes to Fourier spectra
and to continuous-time LFP shot-noise models.

\section{How Bernoulli bins approach a Poisson process}
Partition a fixed duration $T$ into $m=T/\Delta$ bins,
each with independent event probability $p_\Delta=\lambda\Delta$.
Then $N_T\sim\operatorname{Binomial}(m,p_\Delta)$, with
\[
\Prob(N_T=0)=(1-\lambda T/m)^m
\longrightarrow e^{-\lambda T}.
\]
The binomial probability generating function likewise converges
to $\exp[\lambda T(s-1)]$, proving convergence of the
\emph{total count} to a Poisson distribution. Independence
of disjoint bins and consistency across intervals give the
homogeneous Poisson-process limit. A form of Le Cam's inequality
bounds total-variation error in Poisson approximations by a
constant multiple of $\sum_i p_i^2$; hence fixed $T$ and
$p_i=\lambda\Delta$ make this error vanish as $\Delta\to0$.
This approximation is not a guarantee for refractory or bursty
neuronal processes.

\section{Renewal variability}
Suppose interspike intervals $S_1,S_2,\ldots$ are iid,
strictly positive, with mean $m$ and finite variance $v$.
Their coefficient of variation is $CV^2=v/m^2$.
Under standard renewal limit conditions, for long observation
windows $T$ the spike count satisfies
\[
\E[N(T)]\sim T/m,\qquad
\Var[N(T)]\sim (v/m^3)T,\qquad
\frac{\Var N(T)}{\E N(T)}\longrightarrow CV^2.
\]
Thus \emph{long-window} count Fano factor and ISI $CV^2$
are related for a renewal model. Nonrenewal history,
trial mixing, latent gain, short windows and nonstationarity
can break the empirical approximation.

\section{Bartlett spectra and shot-noise LFP}
For a stationary point process, the Bartlett spectrum is the
spectrum of its centered event-counting measure, including
the shot-noise contribution. A homogeneous Poisson train of
rate $\lambda$ has flat continuous spectral density
$S_{nn}(f)=\lambda$ in a two-sided Hz convention
(the DC mean is treated separately).
For a stationary renewal process with inter-event characteristic
function $\varphi(f)=\E[e^{-2\pi ifS}]$, the nonzero-frequency
continuous spectral density takes the form
\[
S_{nn}(f)=\lambda
\frac{1-|\varphi(f)|^2}{|1-\varphi(f)|^2},
\]
under the usual nonlattice and regularity assumptions.
This shows how refractoriness or regular spiking produces
spectral structure without sinusoidal membrane oscillations.

If an idealized LFP component is a linear superposition of
spike-triggered kernels, $y(t)=\int h(t-s)\,dN(s)$,
its centered spectrum satisfies
$S_{yy}(f)=|H(f)|^2S_{nn}(f)$;
independent measurement noise adds its own spectrum.
If both $n$ and $y$ derive from the same events, their cross-spectrum
is proportional to $H(f)S_{nn}(f)$ with sign/conjugation depending
on the cross-spectrum convention. This is an example of
\emph{spike--field coherence without a claim of synaptic causation}.

\section{Hawkes excitation and stability}
For a linear Hawkes intensity
$\lambda(t)=\mu+\sum_{t_j<t}h(t-t_j)$ with
$h\ge0$, $\mu>0$ and branching ratio
$\nu=\int_0^\infty h(s)\,ds<1$,
a stationary version exists under the standard construction,
and its mean rate is $\bar\lambda=\mu/(1-\nu)$.
When $\nu\ge1$, this stationary finite-mean construction
fails in the usual nonnegative linear model; the statement
does not apply without alteration to inhibitory or nonlinear
point-process intensities.

\begin{example}[Arsenal interpretation]\sloppy
Compare \path{spike_count_fano.m},
\path{spike_field_coherence.m} and
\path{spike_field_ppc.m} on Poisson,
renewal, and latent-rate simulations. A function operating on
binned counts or spike phases does not itself implement a
continuous-time Bartlett spectral estimator.
\end{example}